\section{Introduction}
\label{sec:introduction}
% Industrial maintenance and inspection are essential for ensuring the operational continuity and safety of manufacturing facilities.
산업 유지보수 및 점검은 제조 시설의 운영 연속성과 안전을 보장하기 위해 필수적입니다.
% Traditionally, these tasks have relied on human expertise, which is subject to fatigue and inconsistency.
전통적으로 이러한 작업은 인간의 전문성에 의존해 왔으나, 이는 피로도와 일관성 부족에 영향을 받기 쉽습니다.
% With the advent of Industry 4.0, automated inspection systems leveraging deep learning have emerged as a robust alternative. 
인더스트리 4.0의 도래와 함께 딥러닝을 활용한 자동 점검 시스템이 강력한 대안으로 부상했습니다.

% Despite the progress, deploying automated systems in real-world industrial environments remains challenging.
이러한 진전에도 불구하고, 실제 산업 환경에 자동화 시스템을 배치하는 것은 여전히 도전적인 과제입니다.
% Standard computer vision models often struggle with (1) extreme perspective distortions and reflections on gauge surfaces, (2) the dense arrangement of identical indicators leading to matching drift, and (3) the diverse nature of inspection targets that range from simple LEDs to complex semantic signs.
표준 컴퓨터 비전 모델은 (1) 게이지 표면의 극심한 원근 왜곡 및 반사, (2) 동일한 지시등의 밀집 배치로 인한 매칭 드리프트, (3) 단순 LED부터 복잡한 의미론적 표지판에 이르는 점검 대상의 다양성 문제로 인해 어려움을 겪는 경우가 많습니다.
% There is an urgent need for a system that provides high-precision diagnostic capability while remaining computationally efficient for robotic patrolling.
로봇 순찰을 위해 계산 효율성을 유지하면서도 고정밀 진단 능력을 제공하는 시스템에 대한 시급한 요구가 있습니다.

% This paper introduces an integrated diagnostic system that combines state-of-the-art object detection with specialized geometric analysis modules and Vision-Language Models (VLM).
본 논문은 최첨단 객체 탐지 기술과 특화된 기하학적 분석 모듈, 그리고 시각-언어 모델(VLM)을 결합한 통합 진단 시스템을 소개합니다.
% The primary contributions of this work are as follows:
본 연구의 주요 기여도는 다음과 같습니다:
\begin{itemize}
% \item \textbf{Layered Modular Architecture}: We propose a four-layer framework (Orchestration, Perception, Diagnostics, Resource Management) that allows for independent scaling and optimization of each functional unit.
    \item \textbf{계층형 모듈 아키텍처}: 각 기능 단위의 독립적인 확장 및 최적화가 가능한 4계층 프레임워크(오케스트레이션, 인지, 진단, 자원 관리)를 제안합니다.
% \item \textbf{Hybrid Geometric-AI Diagnostics}: We integrate probabilistic deep learning with deterministic geometric refinement (Homography and Hough transforms) to achieve sub-2\% error rates in gauge reading under extreme viewing angles.
    \item \textbf{하이브리드 기하학-AI 진단}: 확률론적 딥러닝과 결정론적 기하학적 정제(호모그래피 및 허프 변환)를 통합하여 극한의 시야각에서도 게이지 판독 오차율 2\% 미만을 달성했습니다.
% \item \textbf{Prioritized Inference Strategy}: We implement a task-specific orchestration logic that selectively invokes intensive VLM reasoning only for complex "ETC" items, utilizing parallel execution to maintain real-time performance (120ms avg. latency).
    \item \textbf{우선순위 기반 추론 전략}: 복잡한 "ETC" 항목에 대해서만 고비용 VLM 추론을 선택적으로 호출하는 작업 특화 오케스트레이션 로직을 구현하고, 병렬 실행을 활용하여 실시간 성능(평균 지연 시간 120ms)을 유지합니다.
% \item \textbf{Robust Empirical Validation}: We demonstrate the system's effectiveness through a comprehensive dataset from four distinct industrial zones, providing detailed qualitative analysis of 12 distinct edge-case scenarios.
    \item \textbf{강건한 실증적 검증}: 4개의 서로 다른 산업 구역에서 수집한 포괄적인 데이터셋을 통해 시스템의 효과를 입증하고, 12가지의 뚜렷한 에지 케이스 시나리오에 대한 상세한 정성적 분석을 제공합니다.
\end{itemize}
